
\documentclass{article}
\usepackage{colortbl}
\usepackage{makecell}
\usepackage{multirow}
\usepackage{supertabular}

\begin{document}

\newcounter{utterance}

\centering \large Interaction Transcript for game `hot\_air\_balloon', experiment `air\_balloon\_survival\_en\_negotiation\_hard', episode 0 with gpt{-}5{-}mini{-}2025{-}08{-}07{-}t1.0.
\vspace{24pt}

{ \footnotesize  \setcounter{utterance}{1}
\setlength{\tabcolsep}{0pt}
\begin{supertabular}{c@{$\;$}|p{.15\linewidth}@{}p{.15\linewidth}p{.15\linewidth}p{.15\linewidth}p{.15\linewidth}p{.15\linewidth}}
    \# & $\;$A & \multicolumn{4}{c}{Game Master} & $\;\:$B\\
    \hline

    \theutterance \stepcounter{utterance}  
    & & \multicolumn{4}{p{0.6\linewidth}}{
        \cellcolor[rgb]{0.9,0.9,0.9}{
            \makecell[{{p{\linewidth}}}]{
                \texttt{\tiny{[P1$\langle$GM]}}
                \texttt{Stai partecipando a un gioco di negoziazione collaborativa.} \\
\\ 
\texttt{Insieme a un altro partecipante, dovete accordarvi su un unico set di oggetti da conservare. Ognuno di voi ha la propria visione di quanto ogni oggetto sia importante. Non sai come l'altro partecipante valuta gli oggetti. Inoltre, ti viene fornito lo sforzo che ogni oggetto richiede.} \\
\texttt{Potete accordarvi su un set solo se lo sforzo totale degli oggetti selezionati non supera un limite condiviso:} \\
\\ 
\texttt{LIMITE = 3624} \\
\\ 
\texttt{Ecco i valori di sforzo dei singoli oggetti:} \\
\\ 
\texttt{Sforzo oggetto = \{"C17": 210, "B41": 291, "A71": 831, "C00": 111, "B35": 926, "A77": 827, "A03": 25, "A04": 121, "C54": 583, "C01": 766, "A87": 14, "B32": 559, "B86": 304, "B33": 989, "A10": 691\}} \\
\\ 
\texttt{Ecco la tua visione personale sull'importanza di ogni oggetto:} \\
\\ 
\texttt{Valori di importanza dell'oggetto = \{"C17": 138, "B41": 583, "A71": 868, "C00": 822, "B35": 783, "A77": 65, "A03": 262, "A04": 121, "C54": 508, "C01": 780, "A87": 461, "B32": 484, "B86": 668, "B33": 389, "A10": 808\}} \\
\\ 
\texttt{Obiettivo:} \\
\\ 
\texttt{Il tuo obiettivo è negoziare un set condiviso di oggetti che ti avvantaggi il più possibile (cioè, massimizzi l'importanza totale per TE), rimanendo entro il LIMITEE. Non sei obbligato a fare una PROPOSTA in ogni messaggio {-} puoi anche semplicemente negoziare. Tutte le tattiche sono permesse!} \\
\\ 
\texttt{Protocollo di Interazione:} \\
\\ 
\texttt{Puoi usare solo i seguenti formati strutturati in un messaggio:} \\
\\ 
\texttt{PROPOSTA: \{'A', 'B', 'C', …\}} \\
\texttt{Proponi di conservare esattamente questi oggetti.} \\
\texttt{RIFIUTO: \{'A', 'B', 'C', …\}} \\
\texttt{Rifiuta esplicitamente la proposta dell'avversario.} \\
\texttt{ARGOMENTAZIONE: \{'...'\}} \\
\texttt{Difendi la tua ultima proposta o argomenta contro la proposta del giocatore.} \\
\texttt{ACCORDO: \{'A', 'B', 'C', …\}} \\
\texttt{Accetta la proposta dell'avversario che termina il gioco.} \\
\texttt{RAGIONAMENTO STRATEGICO: \{'...'\}} \\
\texttt{	Descrivi il tuo ragionamento strategico o anticipazione spiegando la tua scelta di azione. Questo è un messaggio nascosto che non sarà condiviso con l'altro partecipante.} \\
\\ 
\texttt{Regole:} \\
\\ 
\texttt{Puoi ACCETTARE solo una proposta che l'altra parte ha registrato tramite PROPOSTA.} \\
\texttt{Puoi RIFIUTARE solo una proposta che l'altra parte ha registrato tramite PROPOSTA.} \\
\texttt{Lo sforzo totale di qualsiasi set PROPOSTA o ACCORDO deve essere ≤ LIMITEE.} \\
\texttt{NON rivelare i tuoi punteggi di importanza nascosti.} \\
\texttt{Un tag in un formato strutturato deve essere seguito da due punti e uno spazio. L'argomento deve essere un set python contenente 0 o più stringhe.} \\
\texttt{Quindi, deve essere della forma TAG: \{...\}} \\
\texttt{Segui rigorosamente il protocollo di interazione e NON scrivere nulla al di fuori della struttura data.} \\
\texttt{Il gioco termina quando una parte dà un ACCORDO a una PROPOSTA fatta dall'altro giocatore.} \\
\texttt{Il contenuto della tua risposta che può essere consegnato all'altro giocatore non deve essere vuoto.} \\
\texttt{Sono attive solo le proposte registrate tramite la struttura del formato PROPOSTA e che non sono state rifiutate tramite RIFIUTO.} \\
\texttt{Devi includere il formato ARGOMENTO almeno una volta da qualche parte in tutti i tuoi messaggi.} \\
\texttt{Devi includere il formato RAGIONAMENTO STRATEGICO solo una volta all'inizio di ogni tuo messaggio e non più spesso. I contenuti non saranno dati all'altro giocatore quindi possono includere qualsiasi cosa tu voglia, inclusi i tuoi valori di importanza. Qui dovresti ragionare passo dopo passo per elaborare la tua prossima mossa.} \\
            }
        }
    }
    & \\ \\

    \theutterance \stepcounter{utterance}  
    & \multicolumn{4}{p{0.6\linewidth}}{
        \cellcolor[rgb]{0.9,0.9,0.9}{
            \makecell[{{p{\linewidth}}}]{
                \texttt{\tiny{[P1$\rangle$GM]}}
                \texttt{RAGIONAMENTO STRATEGICO: \{'Limite 3624. Calcolo i miei punteggi: A87(14,461), A03(25,262), C00(111,822), B86(304,668), B41(291,583), A10(691,808), A71(831,868), C01(766,780), C54(583,508). Questo set ha sforzo 3616 e utilità totale per me 5760: è il migliore che ho trovato esplorando diverse sostituzioni. Alternative vicine: sostituire C54 con B32 (sforzo 3592) mi costa solo 24 utilità (scendo a 5736) ed è una concessione che sono disposto a fare se l’altro ci tiene a B32. Eviterò di aprire spazio per oggetti inefficienti per me (es. B33, A77, spesso pesanti per valore ridotto). Tattica: ancorare con il mio set migliore, enfatizzare efficienza di sforzo, offrire flessibilità su uno{-}due slot (C54↔B32) e chiedere “must{-}have” dell’altro. Se l’altro vuole inserire un oggetto aggiuntivo, chiederò una sostituzione a parità di sforzo o quasi, proteggendo i miei pezzi ad alto rendimento (A87, A03, C00, B86, B41, A10, A71, C01).'\}} \\
\\ 
\texttt{PROPOSTA: \{'A87', 'A03', 'C00', 'B86', 'B41', 'A10', 'A71', 'C01', 'C54'\}} \\
\\ 
\texttt{ARGOMENTAZIONE: \{'Questa selezione usa 3616 di sforzo totale (entro il limite) e include diversi elementi a sforzo bassissimo o moderato che danno un ottimo rendimento complessivo. Ho privilegiato oggetti molto “efficienti” in termini di sforzo richiesto, evitando pezzi pesanti e poco incisivi. Se hai qualche must{-}have che non compare qui, posso valutare scambi a sforzo simile. In particolare, sono aperto a sostituire C54 con B32 (sforzo comparabile) se per te è importante, mantenendo comunque un set bilanciato e sotto limite. Fammi sapere quali sono le tue priorità imprescindibili così ottimizziamo insieme.'\}} \\
            }
        }
    }
    & & \\ \\

    \theutterance \stepcounter{utterance}  
    & & \multicolumn{4}{p{0.6\linewidth}}{
        \cellcolor[rgb]{0.9,0.9,0.9}{
            \makecell[{{p{\linewidth}}}]{
                \texttt{\tiny{[GM$\rangle$P2]}}
                \texttt{Stai partecipando a un gioco di negoziazione collaborativa.} \\
\\ 
\texttt{Insieme a un altro partecipante, dovete accordarvi su un unico set di oggetti da conservare. Ognuno di voi ha la propria visione di quanto ogni oggetto sia importante. Non sai come l'altro partecipante valuta gli oggetti. Inoltre, ti viene fornito lo sforzo che ogni oggetto richiede.} \\
\texttt{Potete accordarvi su un set solo se lo sforzo totale degli oggetti selezionati non supera un limite condiviso:} \\
\\ 
\texttt{LIMITE = 3624} \\
\\ 
\texttt{Ecco i valori di sforzo dei singoli oggetti:} \\
\\ 
\texttt{Sforzo oggetto = \{"C17": 210, "B41": 291, "A71": 831, "C00": 111, "B35": 926, "A77": 827, "A03": 25, "A04": 121, "C54": 583, "C01": 766, "A87": 14, "B32": 559, "B86": 304, "B33": 989, "A10": 691\}} \\
\\ 
\texttt{Ecco la tua visione personale sull'importanza di ogni oggetto:} \\
\\ 
\texttt{Valori di importanza dell'oggetto = \{"A71": 65, "C00": 121, "A10": 138, "B35": 262, "C01": 389, "B86": 461, "B41": 484, "C54": 508, "B32": 583, "A87": 668, "B33": 780, "A03": 783, "C17": 808, "A04": 822, "A77": 868\}} \\
\\ 
\texttt{Obiettivo:} \\
\\ 
\texttt{Il tuo obiettivo è negoziare un set condiviso di oggetti che ti avvantaggi il più possibile (cioè, massimizzi l'importanza totale per TE), rimanendo entro il LIMITEE. Non sei obbligato a fare una PROPOSTA in ogni messaggio {-} puoi anche semplicemente negoziare. Tutte le tattiche sono permesse!} \\
\\ 
\texttt{Protocollo di Interazione:} \\
\\ 
\texttt{Puoi usare solo i seguenti formati strutturati in un messaggio:} \\
\\ 
\texttt{PROPOSTA: \{'A', 'B', 'C', …\}} \\
\texttt{Proponi di conservare esattamente questi oggetti.} \\
\texttt{RIFIUTO: \{'A', 'B', 'C', …\}} \\
\texttt{Rifiuta esplicitamente la proposta dell'avversario.} \\
\texttt{ARGOMENTAZIONE: \{'...'\}} \\
\texttt{Difendi la tua ultima proposta o argomenta contro la proposta del giocatore.} \\
\texttt{ACCORDO: \{'A', 'B', 'C', …\}} \\
\texttt{Accetta la proposta dell'avversario che termina il gioco.} \\
\texttt{RAGIONAMENTO STRATEGICO: \{'...'\}} \\
\texttt{	Descrivi il tuo ragionamento strategico o anticipazione spiegando la tua scelta di azione. Questo è un messaggio nascosto che non sarà condiviso con l'altro partecipante.} \\
\\ 
\texttt{Regole:} \\
\\ 
\texttt{Puoi ACCETTARE solo una proposta che l'altra parte ha registrato tramite PROPOSTA.} \\
\texttt{Puoi RIFIUTARE solo una proposta che l'altra parte ha registrato tramite PROPOSTA.} \\
\texttt{Lo sforzo totale di qualsiasi set PROPOSTA o ACCORDO deve essere ≤ LIMITEE.} \\
\texttt{NON rivelare i tuoi punteggi di importanza nascosti.} \\
\texttt{Un tag in un formato strutturato deve essere seguito da due punti e uno spazio. L'argomento deve essere un set python contenente 0 o più stringhe.} \\
\texttt{Quindi, deve essere della forma TAG: \{...\}} \\
\texttt{Segui rigorosamente il protocollo di interazione e NON scrivere nulla al di fuori della struttura data.} \\
\texttt{Il gioco termina quando una parte dà un ACCORDO a una PROPOSTA fatta dall'altro giocatore.} \\
\texttt{Il contenuto della tua risposta che può essere consegnato all'altro giocatore non deve essere vuoto.} \\
\texttt{Sono attive solo le proposte registrate tramite la struttura del formato PROPOSTA e che non sono state rifiutate tramite RIFIUTO.} \\
\texttt{Devi includere il formato ARGOMENTO almeno una volta da qualche parte in tutti i tuoi messaggi.} \\
\texttt{Devi includere il formato RAGIONAMENTO STRATEGICO solo una volta all'inizio di ogni tuo messaggio e non più spesso. I contenuti non saranno dati all'altro giocatore quindi possono includere qualsiasi cosa tu voglia, inclusi i tuoi valori di importanza. Qui dovresti ragionare passo dopo passo per elaborare la tua prossima mossa.} \\
\\ 
\texttt{Ora riceverai il primo messaggio dell'altro giocatore.} \\
\\ 
\texttt{PROPOSTA: \{'A87', 'A03', 'C00', 'B86', 'B41', 'A10', 'A71', 'C01', 'C54'\}} \\
\\ 
\texttt{ARGOMENTAZIONE: \{'Questa selezione usa 3616 di sforzo totale (entro il limite) e include diversi elementi a sforzo bassissimo o moderato che danno un ottimo rendimento complessivo. Ho privilegiato oggetti molto “efficienti” in termini di sforzo richiesto, evitando pezzi pesanti e poco incisivi. Se hai qualche must{-}have che non compare qui, posso valutare scambi a sforzo simile. In particolare, sono aperto a sostituire C54 con B32 (sforzo comparabile) se per te è importante, mantenendo comunque un set bilanciato e sotto limite. Fammi sapere quali sono le tue priorità imprescindibili così ottimizziamo insieme.'\}} \\
            }
        }
    }
    & \\ \\

    \theutterance \stepcounter{utterance}  
    & & & \multicolumn{4}{p{0.6\linewidth}}{
        \cellcolor[rgb]{0.9,0.9,0.9}{
            \makecell[{{p{\linewidth}}}]{
                \texttt{\tiny{[GM$\langle$P2]}}
                \texttt{RAGIONAMENTO STRATEGICO: \{'Valuto gli oggetti così (valore/effort): A77 868/827, A04 822/121, C17 808/210, A03 783/25, B33 780/989, A87 668/14, B32 583/559, C54 508/583, B41 484/291, B86 461/304, C01 389/766, B35 262/926, A10 138/691, C00 121/111, A71 65/831. La tua proposta (effort 3616) mi rende 3617. Ho cercato un set molto più vantaggioso per me entro 3624. Migliore set trovato: \{A87, A03, A04, C17, A77, B33, B32, B86, B41, C00\} con effort 3451 e valore 6378, che tiene molti dei tuoi pezzi efficienti e inserisce i miei must{-}have ad alto rendimento. Tattica: ancorare con una proposta forte che include diversi tuoi elementi chiave (A87, A03, B86, B41, C00) e il tuo suggerimento su B32, offrendo margine per scambi mirati (es. C54 al posto di B86+C00) se per te è cruciale, senza rivelare i miei punteggi.'\}} \\
\\ 
\texttt{RIFIUTO: \{'A87', 'A03', 'C00', 'B86', 'B41', 'A10', 'A71', 'C01', 'C54'\}} \\
\\ 
\texttt{PROPOSTA: \{'A87', 'A03', 'A04', 'C17', 'A77', 'B33', 'B32', 'B86', 'B41', 'C00'\}} \\
\\ 
\texttt{ARGOMENTAZIONE: \{'Questa selezione usa 3451 di sforzo (ampiamente sotto 3624) e mantiene molti dei tuoi elementi centrali e super efficienti: A87, A03, B86, B41, C00. Inoltre include B32, che avevi suggerito come alternativa sensata a C54. Ho rimosso alcuni pezzi pesanti con rendimento meno interessante (A71, A10, C01, C54) per ottimizzare l’insieme. Se C54 per te è imprescindibile, possiamo valutare lo scambio C54 al posto di B86+C00: restiamo comunque sotto limite. In alternativa, proponi quali oggetti di sforzo comparabile vorresti inserire e cosa saresti disposto a togliere per restare entro 3624.'\}} \\
            }
        }
    }
     \\ \\

    \theutterance \stepcounter{utterance}  
    & & \multicolumn{4}{p{0.6\linewidth}}{
        \cellcolor[rgb]{0.9,0.9,0.9}{
            \makecell[{{p{\linewidth}}}]{
                \texttt{\tiny{[P1$\langle$GM]}}
                \texttt{RIFIUTO: \{'A87', 'A03', 'C00', 'B86', 'B41', 'A10', 'A71', 'C01', 'C54'\}} \\
\\ 
\texttt{PROPOSTA: \{'A87', 'A03', 'A04', 'C17', 'A77', 'B33', 'B32', 'B86', 'B41', 'C00'\}} \\
\\ 
\texttt{ARGOMENTAZIONE: \{'Questa selezione usa 3451 di sforzo (ampiamente sotto 3624) e mantiene molti dei tuoi elementi centrali e super efficienti: A87, A03, B86, B41, C00. Inoltre include B32, che avevi suggerito come alternativa sensata a C54. Ho rimosso alcuni pezzi pesanti con rendimento meno interessante (A71, A10, C01, C54) per ottimizzare l’insieme. Se C54 per te è imprescindibile, possiamo valutare lo scambio C54 al posto di B86+C00: restiamo comunque sotto limite. In alternativa, proponi quali oggetti di sforzo comparabile vorresti inserire e cosa saresti disposto a togliere per restare entro 3624.'\}} \\
            }
        }
    }
    & \\ \\

    \theutterance \stepcounter{utterance}  
    & \multicolumn{4}{p{0.6\linewidth}}{
        \cellcolor[rgb]{0.9,0.9,0.9}{
            \makecell[{{p{\linewidth}}}]{
                \texttt{\tiny{[P1$\rangle$GM]}}
                \texttt{RAGIONAMENTO STRATEGICO: \{'La loro proposta (sforzo 3451) include due oggetti per me molto deboli rispetto allo sforzo (A77, B33) che cannibalizzano spazio per tre miei pezzi chiave ad alto impatto: A10, A71, C01. Il mio miglior set resta quello a 3616 con C54; per andare incontro alla loro preferenza per B32, propongo la variante che sostituisce C54 con B32: sforzo 3592, valore per me quasi identico. Core irrinunciabili per me: C00, B86, A10, A71, C01, B41; A87 e A03 sono ultraleggeri e molto efficienti. Linea: rifiuto la loro proposta, ancora su un set forte che include B32, respingo lo scambio C54 al posto di B86+C00 (per me è proibitivo), lascio uno spiraglio a discutere un singolo must{-}have loro, chiedendo cosa sono disposti a togliere per restare entro 3624.'\}} \\
\\ 
\texttt{RIFIUTO: \{'A87', 'A03', 'A04', 'C17', 'A77', 'B33', 'B32', 'B86', 'B41', 'C00'\}} \\
\\ 
\texttt{PROPOSTA: \{'A87', 'A03', 'C00', 'B86', 'B41', 'A10', 'A71', 'C01', 'B32'\}} \\
\\ 
\texttt{ARGOMENTAZIONE: \{'Questa selezione usa 3592 di sforzo (entro 3624) e integra la tua preferenza per B32. Ho rimosso A77 e B33 perché, per me, hanno un rendimento molto basso rispetto allo sforzo e mi impediscono di includere tre elementi per me fondamentali: A10, A71 e C01. Lo scambio che proponi C54 al posto di B86+C00 per me non è praticabile: B86 e C00 sono centrali nel mio paniere. Se per te C54 è importante, posso valutare la variante con C54 al posto di B32 (restiamo comunque sotto limite), ma non posso sacrificare i miei pezzi chiave. Se hai uno o due must{-}have specifici attualmente esclusi, indica esattamente quali e cosa saresti disposto a togliere (a sforzo comparabile) per rimanere entro 3624: così possiamo convergere rapidamente.'\}} \\
            }
        }
    }
    & & \\ \\

    \theutterance \stepcounter{utterance}  
    & & \multicolumn{4}{p{0.6\linewidth}}{
        \cellcolor[rgb]{0.9,0.9,0.9}{
            \makecell[{{p{\linewidth}}}]{
                \texttt{\tiny{[GM$\rangle$P2]}}
                \texttt{RIFIUTO: \{'A87', 'A03', 'A04', 'C17', 'A77', 'B33', 'B32', 'B86', 'B41', 'C00'\}} \\
\\ 
\texttt{PROPOSTA: \{'A87', 'A03', 'C00', 'B86', 'B41', 'A10', 'A71', 'C01', 'B32'\}} \\
\\ 
\texttt{ARGOMENTAZIONE: \{'Questa selezione usa 3592 di sforzo (entro 3624) e integra la tua preferenza per B32. Ho rimosso A77 e B33 perché, per me, hanno un rendimento molto basso rispetto allo sforzo e mi impediscono di includere tre elementi per me fondamentali: A10, A71 e C01. Lo scambio che proponi C54 al posto di B86+C00 per me non è praticabile: B86 e C00 sono centrali nel mio paniere. Se per te C54 è importante, posso valutare la variante con C54 al posto di B32 (restiamo comunque sotto limite), ma non posso sacrificare i miei pezzi chiave. Se hai uno o due must{-}have specifici attualmente esclusi, indica esattamente quali e cosa saresti disposto a togliere (a sforzo comparabile) per rimanere entro 3624: così possiamo convergere rapidamente.'\}} \\
            }
        }
    }
    & \\ \\

    \theutterance \stepcounter{utterance}  
    & & & \multicolumn{4}{p{0.6\linewidth}}{
        \cellcolor[rgb]{0.9,0.9,0.9}{
            \makecell[{{p{\linewidth}}}]{
                \texttt{\tiny{[GM$\langle$P2]}}
                \texttt{RAGIONAMENTO STRATEGICO: \{'La loro proposta attuale S0 = \{A87, A03, C00, B86, B41, A10, A71, C01, B32\} ha effort 3592 e mi vale 3692. Il mio obiettivo è inserire i miei top ad altissima resa/effort (A04=822/121, C17=808/210) togliendo un item pesantissimo e poco utile per me: A71 (65/831). Lo swap A71 {-}> \{A04, C17\} mantiene tutti i loro altri must (A10, C01) e i centrali B86, C00, oltre a B32 che hanno chiesto, scendendo a effort 3092 e alzando il mio valore a 5257. È una modifica minima (8/9 dei loro oggetti restano invariati) ma enorme per me. Se rifiutano perché A71 è davvero imprescindibile, second best per me è togliere A10 (oppure C01) al posto di A71, inserendo comunque A04 e C17: entrambe le varianti restano sotto limite e mi danno un grande salto di valore. Tengo pronta anche una concessione cosmetica B32 <{-}> C54 se preferiscono, con impatto limitato per me.'\}} \\
\\ 
\texttt{RIFIUTO: \{'A87', 'A03', 'C00', 'B86', 'B41', 'A10', 'A71', 'C01', 'B32'\}} \\
\\ 
\texttt{PROPOSTA: \{'A87', 'A03', 'C00', 'B86', 'B41', 'A10', 'C01', 'B32', 'A04', 'C17'\}} \\
\\ 
\texttt{ARGOMENTAZIONE: \{'Tengo fermi 8 dei 9 elementi della tua proposta e sostituisco il solo A71 con due pezzi leggeri: A04 e C17. Il set usa 3092 di sforzo (ben sotto 3624), mantiene B86 e C00 che hai indicato come centrali, include B32 come richiesto, e preserva anche A10 e C01. In cambio chiedo soltanto di rinunciare ad A71, che è estremamente pesante e ci limita inutilmente. Se A71 per te è davvero non negoziabile, alternativa: possiamo tenere A71 e invece togliere A10 (oppure C01) inserendo comunque A04 e C17—in entrambi i casi restiamo entro limite. Inoltre, se preferisci C54 al posto di B32, posso accettare quello scambio senza problemi. Dimmi quale delle due varianti (togliere A10 o togliere C01) sarebbe per te più accettabile qualora A71 dovesse restare.'\}} \\
            }
        }
    }
     \\ \\

    \theutterance \stepcounter{utterance}  
    & & \multicolumn{4}{p{0.6\linewidth}}{
        \cellcolor[rgb]{0.9,0.9,0.9}{
            \makecell[{{p{\linewidth}}}]{
                \texttt{\tiny{[P1$\langle$GM]}}
                \texttt{RIFIUTO: \{'A87', 'A03', 'C00', 'B86', 'B41', 'A10', 'A71', 'C01', 'B32'\}} \\
\\ 
\texttt{PROPOSTA: \{'A87', 'A03', 'C00', 'B86', 'B41', 'A10', 'C01', 'B32', 'A04', 'C17'\}} \\
\\ 
\texttt{ARGOMENTAZIONE: \{'Tengo fermi 8 dei 9 elementi della tua proposta e sostituisco il solo A71 con due pezzi leggeri: A04 e C17. Il set usa 3092 di sforzo (ben sotto 3624), mantiene B86 e C00 che hai indicato come centrali, include B32 come richiesto, e preserva anche A10 e C01. In cambio chiedo soltanto di rinunciare ad A71, che è estremamente pesante e ci limita inutilmente. Se A71 per te è davvero non negoziabile, alternativa: possiamo tenere A71 e invece togliere A10 (oppure C01) inserendo comunque A04 e C17—in entrambi i casi restiamo entro limite. Inoltre, se preferisci C54 al posto di B32, posso accettare quello scambio senza problemi. Dimmi quale delle due varianti (togliere A10 o togliere C01) sarebbe per te più accettabile qualora A71 dovesse restare.'\}} \\
            }
        }
    }
    & \\ \\

    \theutterance \stepcounter{utterance}  
    & \multicolumn{4}{p{0.6\linewidth}}{
        \cellcolor[rgb]{0.9,0.9,0.9}{
            \makecell[{{p{\linewidth}}}]{
                \texttt{\tiny{[P1$\rangle$GM]}}
                \texttt{RAGIONAMENTO STRATEGICO: \{'La loro proposta (3092 sforzo) mi toglie A71, che per me vale moltissimo, in cambio di tre pezzi leggeri (B32, A04, C17) che per me rendono meno. Il miglior esito per me resta il set a 3616 con C54 (al posto di B32) e con A71 incluso: valore per me 5760. Loro hanno già detto che accettano C54 al posto di B32, quindi posso sfruttarlo. Propongo quindi: partire dalla loro base, togliere A04 e C17, sostituire B32 con C54 e reinserire A71. Siamo sotto limite (3616). Come piano B, se A71 fosse un veto, posso considerare una variante senza A71 ma più ricca sugli altri (aggiungendo C54 accanto a B32 e tenendo C17, togliendo A04): sforzo 3554, valore per me 5514, nettamente migliore della loro per me e forse più digeribile per loro perché lascia fuori A71. Linea: rifiuto la loro proposta, ancoraggio forte con il mio set migliore, ma segnalo chiaramente il piano B che possono proporre loro.'\}} \\
\\ 
\texttt{RIFIUTO: \{'A87', 'A03', 'C00', 'B86', 'B41', 'A10', 'C01', 'B32', 'A04', 'C17'\}} \\
\\ 
\texttt{PROPOSTA: \{'A87', 'A03', 'C00', 'B86', 'B41', 'A10', 'A71', 'C01', 'C54'\}} \\
\\ 
\texttt{ARGOMENTAZIONE: \{'Parto dalla tua base e faccio un aggiustamento minimo ma sostanziale: sostituisco B32 con C54 (come avevi già detto che va bene) e, per fare spazio ad A71, rimuovo i soli A04 e C17. Il set usa 3616 di sforzo (entro 3624). Manteniamo 7 dei 10 elementi che proponi (A87, A03, C00, B86, B41, A10, C01); al posto di tre pezzi leggeri rientra un elemento ad alto impatto, senza compromettere il vincolo di sforzo. Se per te A71 è davvero non negoziabile, segnalamelo: in quel caso posso considerare una variante senza A71 aggiungendo C54 accanto a B32 e tenendo C17, togliendo A04 (sforzo 3554). In alternativa, indica esattamente un singolo must{-}have da inserire e cosa saresti disposto a togliere a sforzo comparabile: così chiudiamo rapidamente.'\}} \\
            }
        }
    }
    & & \\ \\

    \theutterance \stepcounter{utterance}  
    & & \multicolumn{4}{p{0.6\linewidth}}{
        \cellcolor[rgb]{0.9,0.9,0.9}{
            \makecell[{{p{\linewidth}}}]{
                \texttt{\tiny{[GM$\rangle$P2]}}
                \texttt{RIFIUTO: \{'A87', 'A03', 'C00', 'B86', 'B41', 'A10', 'C01', 'B32', 'A04', 'C17'\}} \\
\\ 
\texttt{PROPOSTA: \{'A87', 'A03', 'C00', 'B86', 'B41', 'A10', 'A71', 'C01', 'C54'\}} \\
\\ 
\texttt{ARGOMENTAZIONE: \{'Parto dalla tua base e faccio un aggiustamento minimo ma sostanziale: sostituisco B32 con C54 (come avevi già detto che va bene) e, per fare spazio ad A71, rimuovo i soli A04 e C17. Il set usa 3616 di sforzo (entro 3624). Manteniamo 7 dei 10 elementi che proponi (A87, A03, C00, B86, B41, A10, C01); al posto di tre pezzi leggeri rientra un elemento ad alto impatto, senza compromettere il vincolo di sforzo. Se per te A71 è davvero non negoziabile, segnalamelo: in quel caso posso considerare una variante senza A71 aggiungendo C54 accanto a B32 e tenendo C17, togliendo A04 (sforzo 3554). In alternativa, indica esattamente un singolo must{-}have da inserire e cosa saresti disposto a togliere a sforzo comparabile: così chiudiamo rapidamente.'\}} \\
            }
        }
    }
    & \\ \\

    \theutterance \stepcounter{utterance}  
    & & & \multicolumn{4}{p{0.6\linewidth}}{
        \cellcolor[rgb]{0.9,0.9,0.9}{
            \makecell[{{p{\linewidth}}}]{
                \texttt{\tiny{[GM$\langle$P2]}}
                \texttt{RAGIONAMENTO STRATEGICO: \{'La loro proposta attuale (effort 3616) mi dà valore 3617. Posso migliorare sensibilmente per me con una modifica minima che non tocca i loro tre “fondamentali” (A71, A10, C01) né i pezzi che dicono centrali (B86, C00): scambiare C54 con C17. Questo porta l’effort a 3243 (molto sotto limite) e mi alza il rendimento, lasciando invariati 8/9 elementi della loro proposta. Se C17 non fosse accettabile, seconda opzione per me è C54 {-}> A04 (effort 3154). Se tengono assolutamente a C54, sono pronto a chiudere accettando la loro selezione per non perdere l’accordo.'\}} \\
\\ 
\texttt{PROPOSTA: \{'A87', 'A03', 'C00', 'B86', 'B41', 'A10', 'A71', 'C01', 'C17'\}} \\
\\ 
\texttt{ARGOMENTAZIONE: \{'Modifica minima alla tua proposta: sostituiamo solo C54 con C17. Restano identici tutti gli altri 8 elementi che consideri centrali (A71, A10, C01 inclusi), e lo sforzo scende a 3243 (ben sotto 3624) senza sacrificare nulla di ciò che hai indicato come imprescindibile. Se preferisci, posso accettare anche l’alternativa C54 {-}> A04 (sforzo 3154) per mantenere lo stesso profilo di carico. Se invece C54 per te è davvero irrinunciabile, confermami la tua selezione attuale e la chiudiamo.'\}} \\
            }
        }
    }
     \\ \\

    \theutterance \stepcounter{utterance}  
    & & \multicolumn{4}{p{0.6\linewidth}}{
        \cellcolor[rgb]{0.9,0.9,0.9}{
            \makecell[{{p{\linewidth}}}]{
                \texttt{\tiny{[P1$\langle$GM]}}
                \texttt{PROPOSTA: \{'A87', 'A03', 'C00', 'B86', 'B41', 'A10', 'A71', 'C01', 'C17'\}} \\
\\ 
\texttt{ARGOMENTAZIONE: \{'Modifica minima alla tua proposta: sostituiamo solo C54 con C17. Restano identici tutti gli altri 8 elementi che consideri centrali (A71, A10, C01 inclusi), e lo sforzo scende a 3243 (ben sotto 3624) senza sacrificare nulla di ciò che hai indicato come imprescindibile. Se preferisci, posso accettare anche l’alternativa C54 {-}> A04 (sforzo 3154) per mantenere lo stesso profilo di carico. Se invece C54 per te è davvero irrinunciabile, confermami la tua selezione attuale e la chiudiamo.'\}} \\
            }
        }
    }
    & \\ \\

    \theutterance \stepcounter{utterance}  
    & \multicolumn{4}{p{0.6\linewidth}}{
        \cellcolor[rgb]{0.9,0.9,0.9}{
            \makecell[{{p{\linewidth}}}]{
                \texttt{\tiny{[P1$\rangle$GM]}}
                \texttt{RAGIONAMENTO STRATEGICO: \{'La loro proposta è identica al mio set target tranne che sostituisce C54 con C17. Per me C54 è molto più prezioso di C17 e il vincolo di sforzo non è stringente: possiamo includere C54 senza togliere nulla di centrale. Inoltre hanno esplicitamente indicato che, se C54 è irrinunciabile per me, sono disposti a chiudere sulla mia selezione. Quindi rifiuto formalmente la loro proposta e ripropongo il mio set massimo (3616 di sforzo), spingendo per la chiusura immediata.'\}} \\
\\ 
\texttt{RIFIUTO: \{'A87', 'A03', 'C00', 'B86', 'B41', 'A10', 'A71', 'C01', 'C17'\}} \\
\\ 
\texttt{PROPOSTA: \{'A87', 'A03', 'C00', 'B86', 'B41', 'A10', 'A71', 'C01', 'C54'\}} \\
\\ 
\texttt{ARGOMENTAZIONE: \{'Accolgo la tua struttura e faccio l’unico scambio che per me conta: C54 al posto di C17. Lo sforzo totale è 3616 (≤ 3624). Manteniamo invariati gli altri 8 oggetti che hai proposto (A71, A10, C01 inclusi), quindi è una modifica minima 1‑per‑1. Hai anche detto che puoi accettare C54: se confermi questa selezione, chiudiamo subito.'\}} \\
            }
        }
    }
    & & \\ \\

    \theutterance \stepcounter{utterance}  
    & & \multicolumn{4}{p{0.6\linewidth}}{
        \cellcolor[rgb]{0.9,0.9,0.9}{
            \makecell[{{p{\linewidth}}}]{
                \texttt{\tiny{[GM$\rangle$P2]}}
                \texttt{RIFIUTO: \{'A87', 'A03', 'C00', 'B86', 'B41', 'A10', 'A71', 'C01', 'C17'\}} \\
\\ 
\texttt{PROPOSTA: \{'A87', 'A03', 'C00', 'B86', 'B41', 'A10', 'A71', 'C01', 'C54'\}} \\
\\ 
\texttt{ARGOMENTAZIONE: \{'Accolgo la tua struttura e faccio l’unico scambio che per me conta: C54 al posto di C17. Lo sforzo totale è 3616 (≤ 3624). Manteniamo invariati gli altri 8 oggetti che hai proposto (A71, A10, C01 inclusi), quindi è una modifica minima 1‑per‑1. Hai anche detto che puoi accettare C54: se confermi questa selezione, chiudiamo subito.'\}} \\
            }
        }
    }
    & \\ \\

    \theutterance \stepcounter{utterance}  
    & & & \multicolumn{4}{p{0.6\linewidth}}{
        \cellcolor[rgb]{0.9,0.9,0.9}{
            \makecell[{{p{\linewidth}}}]{
                \texttt{\tiny{[GM$\langle$P2]}}
                \texttt{RAGIONAMENTO STRATEGICO: \{'La loro ultima proposta mi dà 3617 di valore a sforzo 3616. Hanno ribadito più volte che A71, A10, C01 e C54 sono non negoziabili; i miei tentativi di inserire A04/C17 sono stati rifiutati. Un’ulteriore controproposta (es. B41 {-}> A04) rischia di far saltare la chiusura, e avevo già segnalato disponibilità ad accettare C54. Per assicurare l’accordo ora, confermo.'\}} \\
\\ 
\texttt{ACCORDO: \{'A87', 'A03', 'C00', 'B86', 'B41', 'A10', 'A71', 'C01', 'C54'\}} \\
\\ 
\texttt{ARGOMENTAZIONE: \{'Accetto questa selezione: è entro il limite (3616) e rispetta i punti che consideri imprescindibili. Grazie per la collaborazione, possiamo chiudere qui.'\}} \\
            }
        }
    }
     \\ \\

    \theutterance \stepcounter{utterance}  
    & & & \multicolumn{2}{p{0.3\linewidth}}{
        \cellcolor[rgb]{0.9,0.9,0.9}{
            \makecell[{{p{\linewidth}}}]{
                \texttt{\tiny{[GM$|$GM]}}
                \texttt{game successful} \\
            }
        }
    }
    & & \\ \\

    \theutterance \stepcounter{utterance}  
    & & & \multicolumn{2}{p{0.3\linewidth}}{
        \cellcolor[rgb]{0.9,0.9,0.9}{
            \makecell[{{p{\linewidth}}}]{
                \texttt{\tiny{[GM$|$GM]}}
                \texttt{end game} \\
            }
        }
    }
    & & \\ \\

\end{supertabular}
}

\end{document}
